%----------------------------------------------------------------------------------------
% VAN VRAAG TOT ONDERZOEK. Alle opgaven zijn eigen werk.
%
% Opbouw: Start, leerstof (theorie, werkwijze of uitleg, voorbeeld, korte oefening), en op een
% nieuwe pagina de Opdrachten (door elkaar, stijgende moeilijkheid), Koppel en Test jezelf.
% Labels van dit hoofdstuk beginnen met h1-.
%----------------------------------------------------------------------------------------

\mychapter{1}{Van vraag tot onderzoek}

\lkr{Leerplan B+S'': LPD 55 (onderzoekscompetentie). Deze cursus is geen leerstof die je in één blok geeft: de hoofdstukken ondersteunen een onderzoeksopdracht die over meerdere weken loopt, en kunnen bij elke andere cursus gebruikt worden. Hoofdstuk 1 gaat over de vraag en over wiskundig onderzoek in de enge zin (vermoeden en bewijs), hoofdstuk 2 over modelleren met data, hoofdstuk 3 over rapporteren. De precieze verwachtingen (product, proces, evaluatie) legt de school vast; kijk die na.}

Een onderzoek begint niet met een antwoord, maar met een goede vraag. \textit{Hoe moeilijk is schaken?} is een interessante vraag, maar je kunt ze niet beantwoorden: wat betekent \textit{moeilijk}? Dit hoofdstuk toont hoe je van een onderwerp tot een vraag komt die je met wiskunde kunt beantwoorden, en hoe een wiskundig onderzoek verloopt: experimenteren, een vermoeden opstellen, en het bewijzen of weerleggen. Wat je hier leert, gebruik je in elk onderzoek, ook als het over metingen gaat (hoofdstuk 2).

\begin{oefeningen}[Start]
\oef{h1-start-1}

\oef{h1-verken}
\end{oefeningen}

%========================================================================================
\section{De onderzoeksvraag}

\subsection{De onderzoekscyclus}

Onderzoek doen is iets anders dan een opgave oplossen. Bij een opgave ligt de vraag vast, en weet je dat er een antwoord bestaat. Bij een onderzoek kies je zelf de vraag en de aanpak, en weet je vooraf niet hoe ver je komt. Toch verloopt bijna elk onderzoek in dezelfde fasen, of het nu over getallen, metingen of figuren gaat.

\begin{theorie}
\naast{%
De \textbf{\textcolor{accent}{onderzoekscyclus}} bestaat uit zes fasen:
\begin{itemize}
\item \textbf{oriënteren}: verkennen wat er al bekend is;
\item \textbf{een vraag stellen}: de onderzoeksvraag en deelvragen;
\item \textbf{plannen}: aanpak, middelen en tijd kiezen;
\item \textbf{uitvoeren}: tellen, meten, rekenen, redeneren;
\item \textbf{besluiten}: de vraag beantwoorden;
\item \textbf{rapporteren}: het onderzoek voorstellen en erop terugkijken.
\end{itemize}}{%
\begin{tikzpicture}[font=\scriptsize]
  \foreach \i/\t in {0/oriënteren,1/vraag stellen,2/plannen,3/uitvoeren,4/besluiten,5/rapporteren}{
    \node[draw=accent,fill=white,rounded corners=2pt,inner sep=2pt,minimum width=1.6cm] (n\i) at ({90-60*\i}:1.55) {\t};}
  \foreach \i in {0,...,5}{\pgfmathtruncatemacro{\j}{mod(\i+1,6)}
    \draw[->,accent] (n\i) -- (n\j);}
\end{tikzpicture}}
\end{theorie}

Het is een cyclus, omdat een besluit vaak een nieuwe vraag oproept: dan begint alles opnieuw. Bovendien loop je niet altijd netjes rond. Blijkt bij het uitvoeren dat de vraag te groot is, dan ga je terug en stel je ze bij. Dat is geen mislukking, maar gewoon hoe onderzoek werkt.

\voorbeeld{}{Plaats het onderzoek naar de handdrukken uit oefening~\ref{oef:h1-verken} in de onderzoekscyclus.\label{vb:h1-cyclus}}{%
\uitleg{We gaan na wat er in elke fase gebeurde, ook in de fasen die we bijna zonder nadenken doorliepen.}
\begin{stappenvrij}{2.4cm}{12.3cm}
\text{oriënteren} & Een feest waarop iedereen iedereen een hand geeft: hoeveel handdrukken zijn dat? \\
\text{vraag stellen} & Hoeveel handdrukken zijn er bij $n$ personen? Een getal, dus de vraag is precies. \\
\text{plannen} & Eerst kleine groepen tellen, dan een formule zoeken, dan verklaren. \\
\text{uitvoeren} & Tellen geeft $1, 3, 6, 10$. Vermoeden $\frac{n(n - 1)}{2}$, getest met $6$ personen, en verklaard. \\
\text{besluiten} & Bij $n$ personen zijn er $\frac{n(n - 1)}{2}$ handdrukken. \\
\text{rapporteren} & De formule met haar verklaring opschrijven. Een nieuwe vraag: hoeveel handdrukken zijn er als twee ploegen elkaar na een match de hand geven?
\end{stappenvrij}}

\begin{oefeningen}
\oef{h1-cyclus}
\end{oefeningen}

\subsection{Een goede onderzoeksvraag}

De vraag stuurt het hele onderzoek: ze bepaalt wat je telt, meet of bewijst, en wanneer je klaar bent. Een vage vraag geeft een vaag besluit. Daarom loont het om er tijd in te steken voor je begint.

\begin{theorie}
Een goede \textbf{\textcolor{accent}{onderzoeksvraag}} is
\begin{itemize}
\item \textbf{afgebakend}: ze gaat over één duidelijk omschreven situatie;
\item \textbf{precies}: elk begrip erin is gedefinieerd of meetbaar;
\item \textbf{beantwoordbaar}: met wiskunde, met de middelen en de tijd die je hebt.
\end{itemize}
\end{theorie}

Bij elke eis hoort een typische fout. De tabel toont er telkens een, met een betere vraag ernaast.

\figuren{\renewcommand{\arraystretch}{1.3}%
\begin{tabular}{|>{\centering\arraybackslash}m{2.8cm}|>{\RaggedRight\arraybackslash}m{4.4cm}|>{\RaggedRight\arraybackslash}m{6.6cm}|}
\hline
\textbf{eis} & \textbf{niet goed} & \textbf{beter} \\ \hline
afgebakend & \textit{Hoe werkt het verkeer?} & \textit{Hoe lang moet een verkeerslicht op groen staan opdat 20 wachtende auto's erdoor kunnen?} \\ \hline
precies & \textit{Is een grote pizza voordeliger?} & \textit{Hoe hangt de prijs per cm$^2$ van een pizza af van haar diameter, bij één pizzeria?} \\ \hline
beantwoordbaar & \textit{Hoeveel mensen wonen er in 2200 op aarde?} & \textit{Welk exponentieel model past bij de wereldbevolking van 1950 tot 2020?} \\ \hline
\end{tabular}}

\voorbeeld{}{Maak van het onderwerp \textit{schaken} een goede onderzoeksvraag.\label{vb:h1-vraag}}{%
\uitleg{We vertrekken van een concrete situatie op het bord en zoeken een grootheid die we kunnen tellen of berekenen. Elke poging toetsen we aan de drie eisen.}
\begin{tabular}{@{}>{\RaggedRight\arraybackslash}p{6cm}@{\hspace{4mm}}>{\RaggedRight\small\itshape\arraybackslash}p{8.7cm}@{}}
\textit{Hoe moeilijk is schaken?} & Niet precies: \textit{moeilijk} kun je niet meten. \\[.6ex]
\textit{Hoeveel verschillende schaakpartijen bestaan er?} & Precies, maar niet beantwoordbaar: het aantal is zo groot dat zelfs computers het niet kunnen tellen. \\[.6ex]
\textit{Op hoeveel manieren zet je 8 torens op een schaakbord zo dat ze elkaar niet aanvallen?} & Afgebakend (één bord, 8 torens), precies (een toren valt aan langs zijn rij en zijn kolom) en beantwoordbaar door te tellen. Een goede vraag.
\end{tabular}}

\begin{oefeningen}
\oef{h1-vraag}
\end{oefeningen}

\subsection{Deelvragen}

Ook een goede vraag is vaak te groot om in één keer te beantwoorden. Wie meteen alle opstellingen van 8 torens probeert te tekenen, verdrinkt in de mogelijkheden. Daarom splits je de vraag in stukken die je wel aankunt.

\begin{theorie}
Een \textbf{\textcolor{accent}{deelvraag}} is een kleinere vraag die je apart kunt beantwoorden, en die helpt om de onderzoeksvraag te beantwoorden.
\end{theorie}

\voorbeeld{}{Splits de onderzoeksvraag over de 8 torens uit voorbeeld~\ref{vb:h1-vraag} in deelvragen, en beantwoord ze.\label{vb:h1-deelvragen}}{%
\naast{%
\uitleg{Een bord van $8 \times 8$ is te groot om te tekenen, maar kleine borden lukken wel. We beginnen dus klein, zoeken een patroon en verklaren het.}}{%
\begin{tikzpicture}[scale=.36]
  \foreach \i in {0,...,3}{\foreach \j in {0,...,3}{\pgfmathtruncatemacro{\k}{mod(\i+\j,2)}
    \ifnum\k=0 \fill[accent!18] (\i,\j) rectangle (\i+1,\j+1);\fi}}
  \draw[grijs] (0,0) grid (4,4);
  \foreach \i/\j in {1/0,3/1,0/2,2/3}{\fill[accent] (\i+.5,\j+.5) circle (.28);}
\end{tikzpicture}}
\begin{tabular}{@{}>{\RaggedRight\arraybackslash}p{6.8cm}@{\hspace{4mm}}>{\RaggedRight\small\itshape\arraybackslash}p{7.9cm}@{}}
\textit{Hoeveel opstellingen zijn er op een bord van $1 \times 1$, $2 \times 2$, $3 \times 3$, $4 \times 4$?} & Tekenen en tellen geeft $1$, $2$, $6$, $24$. De figuur toont een van de $24$ opstellingen op een bord van $4 \times 4$. \\[.6ex]
\textit{Welk patroon zit in $1, 2, 6, 24$?} & Het zijn $1!$, $2!$, $3!$ en $4!$. Vermoeden: op een bord van $n \times n$ zijn er $n!$ opstellingen. \\[.6ex]
\textit{Waarom is dat zo?} & In elke rij staat precies één toren. Die van rij 1 heeft $n$ kolommen ter keuze, die van rij 2 nog $n - 1$, enzovoort. Dus $n \cdot (n - 1) \cdots 1 = n!$, en op het schaakbord $8! = 40\,320$. \\[.6ex]
\textit{Wat verandert er met koninginnen?} & Een uitbreiding. Een koningin valt ook schuin aan, en dan bestaat er geen eenvoudige formule meer: op een schaakbord zijn er $92$ opstellingen.
\end{tabular}}

De laatste deelvraag toont dat een uitbreiding veel moeilijker kan zijn dan de hoofdvraag. Daarom komt ze pas op het einde, als de rest af is.

\begin{werkwijze}
Begin met kleine of eenvoudige gevallen. Verander telkens maar één ding (de grootte van het bord, het soort stuk). Zoek dan een patroon, en een verklaring die voor elk geval werkt. Een uitbreiding of veralgemening komt als laatste. Orden de deelvragen zo van eenvoudig naar moeilijk.
\end{werkwijze}

\begin{oefeningen}
\oef{h1-onderwerp}
\end{oefeningen}

%========================================================================================
\section{Experimenteren, vermoeden, bewijzen}

In de wiskunde is een onderzoek pas af als het resultaat bewezen is. De weg ernaartoe is wel experimenteel: je probeert kleine gevallen, zet ze in een tabel, en zoekt een patroon. Dat patroon wordt een vermoeden, en dat vermoeden bewijs je, of je weerlegt het.

\subsection{Een vermoeden opstellen}

\begin{theorie}
Een \textbf{\textcolor{accent}{vermoeden}} of \textbf{\textcolor{accent}{hypothese}} is een uitspraak die je op grond van voorbeelden of metingen verwacht, maar die nog niet bewezen of getoetst is.
\end{theorie}

Een goed vermoeden is zo precies dat je het kunt testen: een formule, of een uitspraak met \textit{altijd} of \textit{nooit}. \textit{Het wordt groter} is te vaag om te testen.

\voorbeeld{}{Hoeveel diagonalen heeft een veelhoek met $n$ hoekpunten? Stel een vermoeden op en test het.\label{vb:h1-diagonalen}}{%
\naast{%
\uitleg{We tekenen en tellen voor kleine $n$, en vergelijken met een tabel die we al kennen: de handdrukken uit oefening~\ref{oef:h1-verken}. Ook daar verbind je elk punt met elk ander punt.}
\begin{stappenvrij}{5.2cm}{5.6cm}
n\colon\ 3,\ 4,\ 5,\ 6 & Kleine veelhoeken tekenen. \\
d\colon\ 0,\ 2,\ 5,\ 9 & Diagonalen tellen, zoals in de zeshoek. \\
h\colon\ 3,\ 6,\ 10,\ 15 & Alle lijnstukken tussen de hoekpunten: de handdrukken $\frac{n(n - 1)}{2}$. \\
h - d\colon\ 3,\ 4,\ 5,\ 6 & Het verschil is telkens $n$: de zijden. \\
d = \frac{n(n - 1)}{2} - n = \frac{n(n - 3)}{2} & Vermoeden, klopt voor alle vier. \\
n = 8\colon\ d = 20 & Test met een nieuw geval: tekenen en tellen geeft ook $20$.
\end{stappenvrij}}{%
\begin{tikzpicture}[scale=1.05]
  \foreach \i in {0,...,5}{\coordinate (P\i) at ({90+60*\i}:1);}
  \draw[accent,thick] (P0) -- (P1) -- (P2) -- (P3) -- (P4) -- (P5) -- cycle;
  \draw[accent!60] (P0) -- (P2) (P0) -- (P3) (P0) -- (P4) (P1) -- (P3) (P1) -- (P4) (P1) -- (P5) (P2) -- (P4) (P2) -- (P5) (P3) -- (P5);
  \foreach \i in {0,...,5}{\fill[accent] (P\i) circle (1.6pt);}
  \node[below,font=\scriptsize] at (0,-1.05) {$n = 6$: $9$ diagonalen};
\end{tikzpicture}}}

Het vermoeden past bij vijf gevallen. Dat maakt het geloofwaardig, maar nog niet zeker: in \ref{sec:h1-bewijzen} bewijzen we het.

\begin{oefeningen}
\oef{h1-vermoeden}
\end{oefeningen}

\subsection{Verschillen in een tabel}

Niet elke tabel lijkt op een tabel die je al kent. Dan helpen de verschillen tussen opeenvolgende waarden. Bij de diagonalen vind je voor $n = 3, \ldots, 7$ de waarden $0, 2, 5, 9, 14$. De verschillen zijn $2, 3, 4, 5$, en de verschillen daarvan zijn telkens $1$. Dat is geen toeval: het wijst op een kwadratische formule. Twee eigenschappen maken dat precies. We schrijven de waarden als een rij $u_n$.

\begin{theorie}
De \textbf{\textcolor{accent}{verschillen}} van een rij zijn $\formule{u_{n+1} - u_n}$. De \textbf{\textcolor{accent}{tweede verschillen}} zijn de verschillen van de verschillen: \quad $\formule{(u_{n+2} - u_{n+1}) - (u_{n+1} - u_n) = u_{n+2} - 2u_{n+1} + u_n}$

\begin{stelling}{Constante verschillen}
Zijn alle verschillen van een rij gelijk aan $d$, dan is \quad $\formule{u_n = dn + q}$ \quad voor een constante $q$.
\end{stelling}

\begin{bewijs}
Stel dat de rij begint bij $n = n_0$. Schrijf $u_n$ als de eerste term plus alle verschillen tot $n$:
\bewijsregel{u_n = u_{n_0} + (u_{n_0+1} - u_{n_0}) + (u_{n_0+2} - u_{n_0+1}) + \cdots + (u_n - u_{n-1}) .}
Dat zijn $n - n_0$ verschillen, elk gelijk aan $d$. Dus
\bewijsregel{u_n = u_{n_0} + (n - n_0)\,d = dn + q \quad\text{met}\quad q = u_{n_0} - n_0 d .}
\end{bewijs}

\begin{stelling}{Constante tweede verschillen}
Zijn alle tweede verschillen van een rij gelijk aan $c$, dan is \quad $\formule{u_n = \tfrac{c}{2}\,n^2 + pn + q}$ \quad voor constanten $p$ en $q$.
\end{stelling}

\begin{bewijs}
Neem $v_n = u_n - \frac{c}{2}n^2$. Omdat $(n + 2)^2 - 2(n + 1)^2 + n^2 = 2$, zijn de tweede verschillen van $v_n$
\bewijsregel{v_{n+2} - 2v_{n+1} + v_n = (u_{n+2} - 2u_{n+1} + u_n) - \tfrac{c}{2} \cdot 2 = c - c = 0 .}
De verschillen van $v_n$ zijn dus allemaal gelijk, zeg aan $p$. Volgens de vorige eigenschap is $v_n = pn + q$, en dus
\bewijsregel{u_n = v_n + \tfrac{c}{2}n^2 = \tfrac{c}{2}n^2 + pn + q .}
\end{bewijs}
\end{theorie}

Constante verschillen betekenen: bij elke stap komt er evenveel bij, en dan is het verband lineair. Constante tweede verschillen betekenen: wat erbij komt, groeit bij elke stap evenveel, en dan is het verband kwadratisch. Omgekeerd heeft elke kwadratische formule constante tweede verschillen, want $c(n + 2)^2 - 2c(n + 1)^2 + cn^2 = 2c$, en een lineair deel geeft tweede verschillen $0$.

Let wel op: de eigenschappen gaan over \textit{alle} verschillen. In een tabel zie je er maar een paar. Dat die constant zijn, is een sterke aanwijzing, maar geen bewijs. In \ref{sec:h1-weerleggen} zie je hoe een patroon na vijf gevallen kan stoppen.

\voorbeeld{}{Met lucifers leg je een vierkant rooster van $n$ op $n$ hokjes. De figuur toont $n = 3$. Zoek een formule voor het aantal lucifers.\label{vb:h1-verschillen}}{%
\naast{%
\uitleg{We tellen voor kleine $n$ en kijken naar de verschillen. Zijn de tweede verschillen constant, dan zoeken we een formule $\frac{c}{2}n^2 + pn + q$, en bepalen we $p$ en $q$ met twee waarden uit de tabel.}
\begin{stappenvrij}{4.9cm}{6.1cm}
n\colon\ 1,\ 2,\ 3,\ 4 & Tekenen en tellen. \\
L\colon\ 4,\ 12,\ 24,\ 40 & \\
\text{verschillen: } 8,\ 12,\ 16 & Niet constant, dus niet lineair. \\
\text{tweede verschillen: } 4,\ 4 & Constant: $c = 4$, dus $L = 2n^2 + pn + q$. \\
2 + p + q = 4 & Invullen van $n = 1$. \\
8 + 2p + q = 12 & Invullen van $n = 2$. \\
p = 2,\quad q = 0 & Tweede vergelijking min de eerste: $6 + p = 8$. \\
L = 2n^2 + 2n & Test met $n = 4$: $32 + 8 = 40$. Klopt.
\end{stappenvrij}}{%
\begin{tikzpicture}[scale=.5]
  \foreach \i in {0,...,3}{\foreach \j in {0,1,2}{
    \draw[accent,line width=1.6pt,shorten <=1.2pt,shorten >=1.2pt] (\j,\i) -- (\j+1,\i);
    \draw[accent,line width=1.6pt,shorten <=1.2pt,shorten >=1.2pt] (\i,\j) -- (\i,\j+1);}}
\end{tikzpicture}}
Het rooster heeft ook een verklaring: er zijn $n + 1$ horizontale lijnen van elk $n$ lucifers, en evenveel verticale. Samen $2 \cdot n(n + 1) = 2n^2 + 2n$, wat het vermoeden bewijst.}

\begin{werkwijze}
Bereken de verschillen. Zijn ze constant, zoek dan $dn + q$. Zijn de tweede verschillen constant, zoek dan $\frac{c}{2}n^2 + pn + q$. Bepaal de onbekenden door twee waarden uit de tabel in te vullen, en test de formule met een waarde die je niet gebruikte. Het resultaat blijft een vermoeden tot je het bewijst.
\end{werkwijze}

\begin{oefeningen}
\oef{h1-verschillen}
\end{oefeningen}

\subsection{Een vermoeden bewijzen}\label{sec:h1-bewijzen}

Het vermoeden over de diagonalen klopt voor $n = 3, 4, 5, 6$ en $8$. Maar klopt het ook voor een veelhoek met $100$ hoekpunten? Dat kun je niet tekenen, en zelfs duizend getekende gevallen garanderen het volgende niet. Alleen een redenering die voor elke $n$ werkt, maakt het zeker.

\begin{theorie}
Een \textbf{\textcolor{accent}{bewijs}} is een redenering die uit definities en eigenschappen die al vaststaan, afleidt dat een uitspraak in alle gevallen waar is. Een bewezen vermoeden heet een \textbf{\textcolor{accent}{stelling}}.
\end{theorie}

\voorbeeld{}{Bewijs dat een veelhoek met $n$ hoekpunten $\frac{n(n - 3)}{2}$ diagonalen heeft.\label{vb:h1-bewijs}}{%
\naast{%
\uitleg{We tellen de diagonalen op een slimme manier: per hoekpunt. Dan moeten we wel opletten dat we niets twee keer tellen.}
\begin{stappenvrij}{3.4cm}{7.6cm}
n - 1 & Een hoekpunt $P$ kun je verbinden met elk ander hoekpunt. \\
n - 3 & Twee daarvan zijn buren van $P$: die lijnstukken zijn zijden. De andere zijn diagonalen. \\
n(n - 3) & Elk van de $n$ hoekpunten heeft $n - 3$ diagonalen. \\
\frac{n(n - 3)}{2} & Zo telden we elke diagonaal twee keer: één keer vanuit elk eindpunt.
\end{stappenvrij}}{%
\begin{tikzpicture}[scale=1.05]
  \foreach \i in {0,...,5}{\coordinate (P\i) at ({90+60*\i}:1);}
  \draw[accent,thick] (P0) -- (P1) -- (P2) -- (P3) -- (P4) -- (P5) -- cycle;
  \draw[accent,very thick,dashed] (P0) -- (P2) (P0) -- (P3) (P0) -- (P4);
  \foreach \i in {0,...,5}{\fill[accent] (P\i) circle (1.6pt);}
  \node[above,font=\scriptsize] at (P0) {$P$};
\end{tikzpicture}}
Elke stap geldt voor elke $n \geqslant 3$, dus het vermoeden is een stelling.}

Merk op dat het bewijs ook \textit{verklaart} waarom de formule klopt. Daarom kost een bewijs zoeken wel moeite, maar levert het ook inzicht op: wie de redenering kent, ziet meteen hoeveel diagonalen er vanuit één hoekpunt vertrekken.

\begin{werkwijze}
\textbf{1.} Bekijk kleine gevallen systematisch en zet de resultaten in een tabel. \textbf{2.} Zoek een patroon (verschillen, verhoudingen, bekende rijen) en formuleer een vermoeden. \textbf{3.} Test het vermoeden op een nieuw geval dat je nog niet gebruikte. \textbf{4.} Zoek een bewijs, of een tegenvoorbeeld. Een patroon in een tabel is nog geen bewijs.
\end{werkwijze}

\begin{oefeningen}
\needspace{10\baselineskip}
\oef{h1-patroon}
\end{oefeningen}

\subsection{Een vermoeden weerleggen}\label{sec:h1-weerleggen}

Een vermoeden kan ook vals zijn. Om dat te tonen is geen lange redenering nodig: één geval waarin het misloopt, volstaat.

\begin{theorie}
Een \textbf{\textcolor{accent}{tegenvoorbeeld}} voor een uitspraak is een geval dat aan alle voorwaarden van de uitspraak voldoet, maar waarin de uitspraak niet geldt. Eén tegenvoorbeeld bewijst dat de uitspraak niet altijd waar is.
\end{theorie}

Hoe bedrieglijk een patroon kan zijn, toont deze bekende puzzel. Zet $n$ punten op een cirkel en verbind ze allemaal met elkaar, zo dat er nooit drie koorden door één punt gaan. Tel de stukken waarin de schijf verdeeld wordt. Voor $n = 1, \ldots, 5$ vind je $1, 2, 4, 8, 16$, en het vermoeden $2^{n-1}$ ligt voor de hand. Maar voor $n = 6$ vind je $31$ stukken, niet $32$.

\figuren{%
\begin{tikzpicture}[scale=.62,font=\scriptsize]
  \draw[grijs] (0,0) circle (1); \fill[accent] (90:1) circle (2pt); \node at (0,-1.45) {$1$};
\end{tikzpicture}\hspace{4mm}
\begin{tikzpicture}[scale=.62,font=\scriptsize]
  \draw[grijs] (0,0) circle (1);
  \foreach \a [count=\i] in {80,250}{\coordinate (P\i) at (\a:1);}
  \draw[accent] (P1) -- (P2); \foreach \i in {1,2}{\fill[accent] (P\i) circle (2pt);} \node at (0,-1.45) {$2$};
\end{tikzpicture}\hspace{4mm}
\begin{tikzpicture}[scale=.62,font=\scriptsize]
  \draw[grijs] (0,0) circle (1);
  \foreach \a [count=\i] in {90,210,330}{\coordinate (P\i) at (\a:1);}
  \foreach \i in {1,...,3}{\foreach \j in {1,...,3}{\ifnum\i<\j \draw[accent] (P\i) -- (P\j);\fi}}
  \foreach \i in {1,...,3}{\fill[accent] (P\i) circle (2pt);} \node at (0,-1.45) {$4$};
\end{tikzpicture}\hspace{4mm}
\begin{tikzpicture}[scale=.62,font=\scriptsize]
  \draw[grijs] (0,0) circle (1);
  \foreach \a [count=\i] in {60,150,220,320}{\coordinate (P\i) at (\a:1);}
  \foreach \i in {1,...,4}{\foreach \j in {1,...,4}{\ifnum\i<\j \draw[accent] (P\i) -- (P\j);\fi}}
  \foreach \i in {1,...,4}{\fill[accent] (P\i) circle (2pt);} \node at (0,-1.45) {$8$};
\end{tikzpicture}\hspace{4mm}
\begin{tikzpicture}[scale=.62,font=\scriptsize]
  \draw[grijs] (0,0) circle (1);
  \foreach \a [count=\i] in {90,162,234,306,18}{\coordinate (P\i) at (\a:1);}
  \foreach \i in {1,...,5}{\foreach \j in {1,...,5}{\ifnum\i<\j \draw[accent] (P\i) -- (P\j);\fi}}
  \foreach \i in {1,...,5}{\fill[accent] (P\i) circle (2pt);} \node at (0,-1.45) {$16$};
\end{tikzpicture}\hspace{4mm}
\begin{tikzpicture}[scale=.62,font=\scriptsize]
  \draw[grijs] (0,0) circle (1);
  \foreach \a [count=\i] in {15,75,130,190,255,320}{\coordinate (P\i) at (\a:1);}
  \foreach \i in {1,...,6}{\foreach \j in {1,...,6}{\ifnum\i<\j \draw[accent] (P\i) -- (P\j);\fi}}
  \foreach \i in {1,...,6}{\fill[accent] (P\i) circle (2pt);} \node at (0,-1.45) {$31$};
\end{tikzpicture}}

Wie hier na vijf gevallen stopt, heeft een fout resultaat. Een patroon testen op een nieuw geval is dus geen formaliteit: het is de eerste kans om een vals vermoeden te betrappen.

\voorbeeld{}{Iemand vermoedt: \textit{voor elk natuurlijk getal $n$ is $n^2 + n + 11$ een priemgetal.} Bewijs het, of geef een tegenvoorbeeld.\label{vb:h1-tegenvoorbeeld}}{%
\uitleg{We testen eerst enkele gevallen. Blijven die kloppen, dan zoeken we gericht een $n$ waarvoor de uitkomst wel een deler heeft.}
\begin{stappenbreed}
n = 0, 1, \ldots, 9\colon\ 11, 13, 17, 23, 31, 41, 53, 67, 83, 101 & Tien priemgetallen na elkaar. Het vermoeden lijkt sterk. \\[1ex]
n^2 + n + 11 = n(n + 1) + 11 & Kies $n$ zo dat ook $n(n + 1)$ een veelvoud van $11$ is: dan is de som deelbaar door $11$. \\[1ex]
n = 10\colon\ 10 \cdot 11 + 11 = 121 = 11^2 & Geen priemgetal: een tegenvoorbeeld.
\end{stappenbreed}
Het vermoeden is vals. Tien kloppende gevallen waren dus geen bewijs, en één geval was genoeg om het te weerleggen.}

\begin{oefeningen}
\oef{h1-tegenvoorbeeld}
\end{oefeningen}

\opdrachtenkop

\begin{oefeningen}
\oef{h1-o8}

\oef{h1-o10}

\oef{h1-o2}

\oef{h1-o9}

\oef{h1-o1}

\oef{h1-o11}

\oef{h1-o3}

\oef{h1-o16}

\oef{h1-o13}

\oef{h1-o4}

\oef{h1-o14}

\oef{h1-o12}

\oef{h1-o5}

\oef{h1-o18}

\oef{h1-o15}

\oef{h1-o17}

\oef{h1-o6}

\oef{h1-o20}

\oef{h1-o19}

\oef{h1-o7}
\end{oefeningen}

%----------------------------------------------------------------------------------------
% KOPPEL: verbanden met andere leerstof
%----------------------------------------------------------------------------------------
\needspace{20\baselineskip}
\opdrachtsectie{Koppel}

\lkr{Logica en bewijzen, Telproblemen en Rijen: kijk na wat al gezien is. Het pakket legt de volgorde niet vast.}

\begin{oefeningen}
\oef{h1-koppel-1}

\oef{h1-koppel-2}

\oef{h1-koppel-3}
\end{oefeningen}

%----------------------------------------------------------------------------------------
\opdrachtsectie{Test jezelf}

\begin{oefeningen}
\oef{h1-t1}

\oef{h1-t2}

\oef{h1-t3}

\oef{h1-t4}

\oef{h1-t5}
\end{oefeningen}

\herhaalschema{%
  \fv{oef:h1-t1}{vb:h1-vraag}{oef:h1-vraag}%
  \fv{oef:h1-t2}{vb:h1-diagonalen}{oef:h1-vermoeden}%
  \fv{oef:h1-t3}{vb:h1-tegenvoorbeeld}{oef:h1-tegenvoorbeeld}%
  \fv{oef:h1-t4}{vb:h1-deelvragen}{oef:h1-onderwerp}%
  \fv{oef:h1-t5}{vb:h1-bewijs}{oef:h1-patroon}}

\ifzelfstudie
\section*{Maak je eigen samenvatting}
Sluit dit boek en schrijf in je schrift: de fasen van de onderzoekscyclus, de drie eisen aan een goede onderzoeksvraag, wat een deelvraag, een vermoeden, een bewijs en een tegenvoorbeeld is, de twee eigenschappen over verschillen (met hun bewijs), en de vier stappen van een wiskundig onderzoek, met een voorbeeld van een vermoeden dat fout bleek. Vergelijk daarna met de kaders.
\fi
